\begin{example}[A field and a polynomial ring]
Every field $K$ is reduced.  More generally, every integral domain is reduced, so
\[
k[t],\qquad k[x,y],\qquad \ZZ
\]
are reduced.  Reducedness is weaker than being a domain: later examples exhibit reduced
rings with zero divisors.
\end{example}

\begin{example}[Squarefree versus nonsquarefree integers]
The ring $\ZZ/6\ZZ$ is reduced because
\[
(6)=(2)\cap(3)
\]
is a radical ideal of $\ZZ$.  By contrast $\ZZ/12\ZZ$ is not reduced: the nonzero class
$\bar 6$ satisfies
\[
\bar6^2=\overline{36}=0.
\]
More generally, $\ZZ/n\ZZ$ is reduced exactly when $n$ is squarefree.
\end{example}

\begin{example}[Dual numbers]
Let
\[
A=k[\varepsilon]/(\varepsilon^2).
\]
Then $\bar\varepsilon\neq0$ but $\bar\varepsilon^2=0$, so $A$ is nonreduced.  Its
nilradical is $(\bar\varepsilon)$ and
\[
A_{\mathrm{red}}\cong k.
\]
Both spectra consist of one point.
\end{example}

\begin{example}[Higher doubled points]
For
\[
A_n=k[t]/(t^n),\qquad n\ge2,
\]
the class $\bar t$ is nilpotent of order $n$.  The only prime ideal is $(\bar t)$, so all
$A_n$ have the same one-point spectrum and the same reduction $k$.  Their rings are
nonisomorphic for different $n$ because the maximum nilpotency order changes.
\end{example}

\begin{example}[A double line]
Set
\[
A=k[x,y]/(y^2).
\]
The nilradical is $(\bar y)$ and
\[
A_{\mathrm{red}}
\cong
k[x,y]/(y)
\cong
k[x].
\]
Thus the double line and the ordinary line $y=0$ have identical prime spectra but different
coordinate rings.
\end{example}

\begin{example}[Crossing axes: reduced but reducible]
The ring
\[
A=k[x,y]/(xy)
\]
is reduced because $(xy)$ is radical: in the UFD $k[x,y]$, the element $xy$ is squarefree.
But $A$ is not a domain because
\[
\bar x\bar y=0
\]
with $\bar x,\bar y\neq0$.  Correspondingly,
\[
\Spec A=V(\bar x)\cup V(\bar y)
\]
is reducible.
\end{example}

\begin{example}[Irreducible but nonreduced]
Let
\[
A=k[x,y]/(x^2).
\]
Then $\Nil(A)=(\bar x)$ and
\[
A_{\mathrm{red}}\cong k[y].
\]
Hence $\Spec A$ is homeomorphic to the irreducible affine line.  Nevertheless $A$ is
nonreduced.  Its generic point is $(\bar x)$, not $(0)$.
\end{example}

\begin{example}[Computing a reduction]
Let
\[
A=k[x,y]/(x^2,xy).
\]
The ideal $(x^2,xy)$ has radical $(x)$, because every common zero of $x^2$ and $xy$ lies on
$x=0$, and algebraically $x^2\in I$ implies $x\in\sqrt I$.  Therefore
\[
A_{\mathrm{red}}\cong k[x,y]/(x)\cong k[y].
\]
The nilpotent class $\bar x$ is killed by reduction.
\end{example}

\begin{example}[A reduced product]
The ring
\[
k\times k
\]
is reduced because each factor is reduced.  It has zero divisors and a disconnected
spectrum consisting of two points.  Thus reducedness does not imply connectedness or
irreducibility.
\end{example}

\begin{example}[A quotient of a reduced ring can be nonreduced]
The polynomial ring $k[t]$ is reduced, but
\[
k[t]/(t^2)
\]
is not.  The issue is not the ambient ring; it is that the ideal $(t^2)$ is not radical.
Replacing it by its radical $(t)$ produces the reduced quotient.
\end{example}

\begin{example}[Localization can remove nilpotent structure]
Let
\[
A=k[x,y]/(x^2,xy).
\]
The ring is nonreduced because $\bar x^2=0$.  On the principal open $D(\bar y)$, however,
$\bar y$ becomes invertible and the relation $\bar x\bar y=0$ forces $\bar x=0$.  Thus
\[
A_{\bar y}\cong k[y,y^{-1}],
\]
which is reduced.  Nilpotent structure may therefore be supported on a proper closed subset.
\end{example}

\begin{example}[Reduction commutes with localization]
For the same ring
\[
A=k[x,y]/(x^2,xy),
\]
we have $A_{\mathrm{red}}\cong k[y]$.  Localizing at $y$ gives
\[
(A_{\mathrm{red}})_y\cong k[y,y^{-1}].
\]
On the other hand $A_y\cong k[y,y^{-1}]$ is already reduced.  Hence in this example
\[
(A_y)_{\mathrm{red}}\cong(A_{\mathrm{red}})_y.
\]
\end{example}

\begin{example}[A finite reduced algebra]
Assume $k$ is algebraically closed.  The quotient
\[
A=k[t]/((t-a)(t-b)(t-c))
\]
with $a,b,c$ distinct is reduced.  The Chinese remainder theorem gives
\[
A\cong k\times k\times k.
\]
Geometrically this is a finite set of three reduced points.
\end{example}

\begin{example}[Reducedness can fail after field extension]
Let $k=\FF_p(s)$ and
\[
A=k[x]/(x^p-s).
\]
The polynomial $x^p-s$ is irreducible over $k$, so $A$ is a field.  After adjoining
$\alpha=s^{1/p}$,
\[
A\otimes_k k(\alpha)
\cong
k(\alpha)[x]/((x-\alpha)^p),
\]
which is nonreduced.  This is the basic warning behind geometric reducedness.
\end{example}
